\documentclass[12pt,a4paper]{article}

% ─── Packages ────────────────────────────────────────────────────────────────
\usepackage[a4paper, margin=2.5cm]{geometry}
\usepackage[utf8]{inputenc}
\usepackage[T1]{fontenc}
\usepackage{lmodern}
\usepackage{microtype}
\usepackage{setspace}
\usepackage{titlesec}
\usepackage{amsmath,amssymb}
\usepackage{booktabs}
\usepackage{longtable}
\usepackage{array}
\usepackage{graphicx}
\usepackage{xcolor}
\usepackage{hyperref}
\usepackage[square,numbers,sort&compress]{natbib}
\usepackage{caption}
\usepackage{enumitem}
\usepackage{fancyhdr}

% ─── Layout ──────────────────────────────────────────────────────────────────
\onehalfspacing
\setlength{\parindent}{1.5em}
\setlength{\parskip}{0.4em}

\titleformat{\section}{\large\bfseries}{}{0em}{\thesection\quad}
\titleformat{\subsection}{\normalsize\bfseries}{}{0em}{\thesubsection\quad}

\hypersetup{
  colorlinks=true,
  linkcolor=blue!60!black,
  citecolor=blue!60!black,
  urlcolor=blue!60!black,
  pdftitle={EcoSynth: A Neuro-Symbolic Framework for Green Retrosynthesis Planning},
  pdfauthor={Jeyadev}
}

\pagestyle{fancy}
\fancyhf{}
\rhead{\small EcoSynth: Neuro-Symbolic Green Retrosynthesis}
\lhead{\small \thepage}
\renewcommand{\headrulewidth}{0.4pt}

% ─── Title Block ─────────────────────────────────────────────────────────────
\begin{document}

\begin{titlepage}
  \centering
  \vspace*{2cm}
  {\LARGE\bfseries EcoSynth: A Neuro-Symbolic Framework for\\[0.5em]
  Green Retrosynthesis Planning with\\[0.5em]
  Dynamic Constraint Propagation and\\[0.5em]
  Hallucination-Aware Route Validation\par}
  \vspace{2cm}
  {\large Jeyadev\par}
  \vspace{0.5cm}
  {\normalsize Department of Computer Science and Engineering\par}
  \vspace{0.3cm}
  {\normalsize \today\par}
  \vspace{2cm}
  \hrule
  \vspace{0.5cm}
  \begin{abstract}
    \noindent
    Retrosynthesis—the decomposition of a target molecule into commercially
    available precursors—has long been a central challenge in organic chemistry.
    Recent deep learning methods have substantially improved prediction accuracy,
    yet they offer little transparency into why a particular route is chosen, and
    they rarely account for environmental sustainability. In this work, I propose
    EcoSynth, a neuro-symbolic system that integrates retrieval-augmented
    generation (RAG), large language model (LLM) candidate generation, dynamic
    constraint propagation, and a multi-tier hallucination taxonomy to produce
    chemically valid, interpretable, and green retrosynthetic routes. The system
    retrieves precedent reactions from a corpus of 485{,}753 USPTO reactions
    using sentence embeddings, propagates reaction-type constraints
    monotonically through a formal constraint graph $G_C^{(k)}$, and scores
    each candidate route using a composite metric
    $\textit{RouteScore} = \alpha G + \beta V + \gamma C + \delta H$ that
    explicitly balances greenness, validity, confidence, and human preference.
    A five-class hallucination taxonomy (HT-01 through HT-05) flags valence
    errors, stereo violations, unsupported precedents, reagent incompatibilities,
    and constraint violations at the node level. Empirical evaluation on a
    set of benchmark molecules demonstrates that EcoSynth consistently returns
    multi-step routes with higher greenness scores than unconstrained baselines,
    while the reasoning timeline component makes each decision step auditable.
    The work exposes several genuine limitations—including the inadequacy of
    small LLMs for retrosynthesis generation and the computational cost of
    dense vector retrieval at scale—and suggests concrete directions for
    improvement.
  \end{abstract}
  \vspace{0.5cm}
  \hrule
  \vspace{0.5cm}
  {\small \textbf{Keywords:} retrosynthesis, green chemistry, retrieval-augmented generation,
  constraint propagation, hallucination detection, neuro-symbolic AI, LLM chemistry}
\end{titlepage}

\tableofcontents
\newpage

% ─── 1. Introduction ─────────────────────────────────────────────────────────
\section{Introduction}

The design of synthetic routes for complex organic molecules is one of the most
cognitively demanding tasks in chemistry. For decades it relied almost entirely
on human expertise: a trained chemist would mentally apply retrosynthetic
disconnection rules, consult reaction databases, and draw on years of
laboratory experience to trace a path from a target structure back to simple,
affordable starting materials. Corey and Wipke's seminal 1969 work
\cite{corey1969lhasa} was the first serious attempt to formalise this process
computationally, encoding expert heuristics into the LHASA program. That effort
established a tradition of rule-based, symbolic planning systems that remained
dominant through the 1990s and into the early 2000s.

The landscape shifted markedly with the application of deep learning to
chemistry. Segler et al.~\cite{segler2018mcts} demonstrated that a Monte Carlo
tree search guided by neural networks could match the performance of expert
chemists on multi-step retrosynthesis for a curated set of molecules—a result
that generated considerable excitement and some controversy regarding evaluation
methodology. Around the same time, sequence-to-sequence transformers trained on
USPTO reaction data began to show remarkable accuracy in predicting reaction
products and, when run in reverse, in suggesting retrosynthetic disconnections
\cite{schwaller2019molecular,liu2017retrosynthesis}. The field rapidly
bifurcated into template-based approaches (which apply extracted reaction
templates to transform a target structure) and template-free methods (which
treat the problem as a sequence translation task).

Both branches have achieved impressive benchmark results, yet both share a
significant blind spot: they optimise for chemical feasibility and reaction
yield while paying little attention to environmental impact. The green chemistry
movement, formalised by Anastas and Warner's twelve principles \cite{anastas1998greenchemistry},
provides a well-established framework for evaluating the environmental cost of
a synthesis route—through metrics such as atom economy, E-factor, and process
mass intensity (PMI). However, these metrics are rarely integrated into
automated synthesis planning tools in any rigorous way. A second, related
problem is interpretability. When a neural model suggests a retrosynthetic
disconnection, it offers no explanation beyond the output itself. For a
practising chemist, understanding why a particular bond was cleaved, which
database precedent supports the proposed reaction, and whether any constraint
was violated in the process is as important as the suggestion itself.

In this work I address both gaps by developing EcoSynth, a system that combines
the complementary strengths of neural retrieval and symbolic reasoning. The
core contributions are as follows. First, I propose a dynamic constraint graph
$G_C^{(k)}$ that monotonically restricts the space of allowable reaction types,
solvents, and reagents as each intermediate in a route is accepted. Second, I
introduce a five-class hallucination taxonomy that catches distinct failure
modes at the symbolic level—before they propagate into a route score. Third,
I embed green chemistry metrics directly into the route scoring function,
so that sustainability is not an afterthought but a primary objective. Fourth,
I show that a reasoning timeline, exposing every RAG retrieval, constraint
update, and firewall decision, is technically tractable and practically useful.

The remainder of this report is organised as follows. Section~2 surveys the
relevant literature across retrosynthesis, green chemistry, RAG, and
hallucination detection. Section~3 identifies the specific research gap that
motivates EcoSynth. Section~4 states the problem formally. Section~5 describes
the system architecture and its components in detail. Sections~6 and~7 present
the analysis and experimental results. Section~8 discusses the findings
critically. Sections~9 and~10 address limitations and future work.
Section~11 concludes.

% ─── 2. Literature Review ────────────────────────────────────────────────────
\section{Literature Review}

\subsection{Automated Retrosynthesis: Foundational Work}

The foundational era of computer-assisted retrosynthesis began with Corey and
Wipke \cite{corey1969lhasa}, whose LHASA system encoded disconnection
heuristics as production rules applied to a graph representation of the target
molecule. Wipke and Howe \cite{wipke1977secs} extended this with the SECS
system, adding more sophisticated atom mapping. These programs worked within
narrow structural domains and required expert curation of the rule base; they
could not generalise to chemistries not represented in their hand-crafted
knowledge base. Szymku{\'c} et al.~\cite{szymkuc2016casp} provide a thorough
retrospective of this era, noting that the systems of the 1980s and 1990s
were ``brittle by design''—their precision came at the cost of recall, and
routes outside the rule base were simply invisible to them.

The transition to learned models began with the availability of large reaction
databases, particularly the USPTO corpus assembled by Lowe
\cite{lowe2012uspto}. Coley et al.~\cite{coley2017synthesis} were among the
first to demonstrate that molecular similarity to known reactions could serve as
a surprisingly effective retrosynthetic heuristic—an approach that anticipated
the later success of retrieval-augmented methods. Liu et
al.~\cite{liu2017retrosynthesis} showed that a sequence-to-sequence recurrent
network trained on reaction SMILES could reproduce known retrosynthetic
disconnections with around 55\% top-1 accuracy on the USPTO-50K benchmark,
establishing this dataset as a standard evaluation arena.

\subsection{Deep Learning for Retrosynthesis}

The transformer architecture \cite{vaswani2017attention} was quickly adapted to
chemistry. Schwaller et al.~\cite{schwaller2019molecular} trained the Molecular
Transformer on 480K USPTO reactions and reported top-1 accuracy of 88.7\% on
reaction prediction, substantially higher than earlier RNN approaches. Applied
in retrosynthesis mode, the same architecture reached competitive
single-step accuracy, though multi-step planning remained an open challenge.

Segler et al.~\cite{segler2018mcts} tackled multi-step planning through a
neural-guided Monte Carlo tree search, combining a learned policy network with
a value network to prioritise promising search paths. Their 2018 Nature paper
reported that the system could solve 60\% of a set of complex natural product
targets, with routes rated as highly as those proposed by expert chemists in a
double-blind evaluation. The paper generated methodological debate—subsequent
work questioned whether the evaluation set was truly unbiased and whether
matching chemist approval is the right metric for a system that might propose
routes along dimensions (cost, greenness, throughput) that experts historically
have not prioritised.

Chen et al.~\cite{chen2020retro} introduced Retro*, an AND-OR graph search
guided by a learned value function, which achieved state-of-the-art
single-step coverage on USPTO-50K. Genheden et al.~\cite{genheden2020aizynthfinder}
released AizynthFinder, an open-source MCTS implementation built on top of
RDKit \cite{landrum2006rdkit} that has become the \emph{de facto} benchmarking
platform for multi-step planning. EcoSynth integrates AizynthFinder as its
planning backbone, though the current prototype falls back to RAG-based
single-step routes when the MCTS search fails to find a solution within its
time budget—an honest limitation discussed further in Section~9.

Graph neural network formulations have become increasingly popular.
Dai et al.~\cite{dai2019retrosynthesis} used a conditional graph logic network
that models retrosynthetic rules as probabilistic logical forms. Somnath
et al.~\cite{somnath2021learning} proposed a graph edit distance formulation
that avoids the need for atom mapping. Yan et al.~\cite{yan2022retrocomposer}
showed that composing multiple partial templates can recover valid routes that
individual templates miss. A common limitation across these approaches is that
they optimise reconstruction accuracy on held-out USPTO reactions, which may
not reflect the distribution of molecules encountered in drug discovery or
process chemistry.

\subsection{Green Chemistry and Computational Metrics}

The twelve principles of green chemistry \cite{anastas1998greenchemistry}
provide a conceptual framework that has been widely adopted in pharmaceutical
and fine chemical manufacturing. Of the quantitative metrics derived from these
principles, atom economy \cite{sheldon2007efactor} and the E-factor (mass of
waste per mass of product) are the most commonly cited. Sheldon's 2007 review
\cite{sheldon2007efactor} surveys their application across industrial sectors
and notes that fine chemicals and pharmaceuticals consistently show the worst
E-factors—a finding that motivates automated optimisation.

Constable et al.~\cite{constable2007solvents} demonstrated that solvent
selection is the single largest driver of E-factor in typical pharmaceutical
syntheses, accounting for 80--90\% of total mass of materials used.
Henderson et al.~\cite{henderson2011chem21} formalised this insight into the
CHEM21 solvent selection guide, which classifies solvents on a hazard and
environmental scorecard. EcoSynth incorporates the CHEM21 classification
directly into its green scorer. Process mass intensity (PMI), proposed by
Roschangar et al.~\cite{roschangar2015pmi}, extends atom economy to account
for all inputs including solvents and reagents; it has been adopted as an
industry standard for process chemistry.

One notable gap in the literature is the near-complete absence of green metrics
in automated retrosynthesis planners. Mikulak-Klucznik et al.'s 2020 Nature
paper \cite{mikulak2020computational} on Chematica's planning of natural
product syntheses mentions sustainability only in passing, and the criteria
used to select among competing routes are dominated by synthetic accessibility
and step count. Struble et al.~\cite{struble2020current} survey the current
state of AI in medicinal chemistry and note that ``green and sustainable
synthesis remains an underexplored direction for computational tools.'' This
observation directly motivates the design of EcoSynth's scoring function.

\subsection{Retrieval-Augmented Generation}

Retrieval-augmented generation, introduced in the NLP community by Lewis et
al.~\cite{lewis2020rag}, addresses the tendency of language models to hallucinate
facts by grounding generation in retrieved documents. The architecture—a
retriever that fetches relevant context followed by a generator that conditions
on that context—has since found application far beyond text generation.
Reimers and Gurevych~\cite{reimers2019sentencebert} introduced
Sentence-BERT, which produces dense semantic embeddings for natural language
sentences that can be indexed for efficient similarity search. Johnson et
al.~\cite{johnson2019faiss} developed FAISS, the approximate nearest-neighbour
library that underlies most production vector search systems. EcoSynth uses
ChromaDB as a persistent vector store, with all-MiniLM-L6-v2 as the embedding
model—a lightweight but effective variant of the sentence transformer family.

In chemistry, Schwaller et al.'s RXNFP \cite{schwaller2021rxnfp} showed that
reaction fingerprints derived from an attention-based encoder map reactions into
a smooth continuous space where chemically similar transformations cluster
together. This opened the door to similarity-based retrieval as a principled
approach to finding relevant reaction precedents. EcoSynth's retrieval module
builds on this insight, though it uses text-based embeddings of reaction
metadata rather than specialised reaction fingerprints—a practical compromise
driven by the availability of pre-trained general-purpose encoders.

The application of LLMs to chemistry has accelerated rapidly. Bran et
al.~\cite{bran2023chemcrow} showed that a GPT-4 agent equipped with chemistry
tools (including RDKit functions, reaction databases, and safety checkers) could
complete multi-step chemistry research tasks. Boiko et al.~\cite{boiko2023autonomous}
demonstrated a GPT-4 system that autonomously conducts chemical experiments by
interfacing with laboratory robots. White et al.~\cite{white2023chemllm}
assessed several LLMs on chemistry benchmarks and found significant variation
in performance—with smaller models often failing on tasks requiring precise
structural reasoning. This last finding is directly relevant to EcoSynth, as
the 1.5-billion-parameter Gemma-3-1B model used for candidate generation proved
insufficient for reliable SMILES generation, an issue discussed candidly in
Section~9.

\subsection{Hallucination in Language Models}

Ji et al.~\cite{ji2023hallucination} provide the most comprehensive taxonomy of
hallucination in NLG systems, distinguishing between intrinsic hallucinations
(contradicting source material) and extrinsic hallucinations (adding
unsupported information). Huang et al.~\cite{huang2023hallucination} survey
mitigation strategies, noting that retrieval augmentation reduces but does not
eliminate hallucination. In the chemistry context, the consequences of
hallucination are particularly severe: a proposed SMILES string that fails RDKit
sanitisation corresponds to a structure that cannot exist, while a plausible but
unsupported reaction may waste expensive materials or produce hazardous by-products.

Krenn et al.~\cite{krenn2020selfies} developed SELFIES, a string representation
that guarantees syntactic validity for any token sequence—addressing HT-01 at
the representational level. However, syntactic validity does not imply chemical
plausibility; a valid SMILES may still violate valence rules, contain
stereochemically impossible configurations, or describe a molecule with no known
reactivity. EcoSynth's firewall addresses this by applying a hierarchy of checks
from syntactic (RDKit parse) through learned plausibility (ChemBERTa-77M
classifier \cite{chithrananda2020chemberta}) to database support (HT-03 precedent
check). The approach is analogous to a multi-level type system: catching errors
at the earliest possible stage reduces computational waste.

\subsection{Molecular Machine Learning}

ChemBERTa \cite{chithrananda2020chemberta} adapts BERT-style pretraining
\cite{devlin2019bert} to molecular SMILES strings, producing representations
that capture local chemical context. Ross et al.~\cite{ross2022chemberta2}
extended this to a larger pretraining corpus and demonstrated improvements on
several molecular property benchmarks. EcoSynth uses ChemBERTa-77M as a binary
classifier distinguishing valid from invalid reaction intermediates.

Schneider et al.~\cite{schneider2016classification} developed a 50-class
reaction classification scheme for USPTO reactions that underpins much of the
supervised learning work in retrosynthesis. The classification, which includes
classes such as \textit{buchwald\_hartwig}, \textit{suzuki\_coupling},
\textit{diels\_alder}, and \textit{reductive\_amination}, is used in EcoSynth
both as a metadata field for retrieved precedents and as a component of the
constraint graph's allowed reaction type set $T$.

% ─── 3. Research Gap ─────────────────────────────────────────────────────────
\section{Research Gap}

Three specific gaps motivate this work.

\textbf{Gap 1: Sustainability is absent from most planners.}
Despite the maturity of green chemistry metrics and their clear industrial
relevance, I am not aware of any publicly available retrosynthesis system that
integrates atom economy, E-factor, PMI, and solvent hazard into a unified
scoring function that directly influences route selection. Existing systems—
including AizynthFinder \cite{genheden2020aizynthfinder}, Retro* \cite{chen2020retro},
and the Molecular Transformer \cite{schwaller2019molecular}—optimise for
chemical feasibility or step count and treat green metrics, if at all, as
post-hoc annotations.

\textbf{Gap 2: Hallucinations from LLMs are not systematically categorised.}
The use of LLMs as components of chemistry pipelines is growing, but the
failure modes specific to molecular generation have not been catalogued with
the same rigour as hallucinations in text generation. The five-class HT
taxonomy introduced here—valence violation, stereo impossibility, no precedent,
reagent incompatibility, constraint violation—represents an initial attempt at
such a domain-specific categorisation.

\textbf{Gap 3: Reasoning is opaque.}
Even systems that offer explanations typically do so at the route level (``this
route was chosen because it has the highest score'') rather than at the
decision level (``this intermediate was accepted because its precedent
similarity was 0.73 and its CHEM21 solvent score is A; constraint propagation
then removed 12 reaction types from the constraint graph''). The reasoning
timeline in EcoSynth is a response to this gap.

% ─── 4. Problem Statement ─────────────────────────────────────────────────────
\section{Problem Statement}

Given a target molecule $M_{\text{target}}$ specified as a SMILES string, a
user-supplied preference vector $\mathbf{p} = (\alpha, \beta, \gamma, \delta)$
with $\sum p_i = 1$, and an optional set of pinned intermediates
$\mathcal{P} \subseteq \mathcal{M}$, the task is to return a ranked list of
retrosynthetic routes $\{r_1, r_2, \ldots, r_K\}$ such that:

\begin{enumerate}[label=(\roman*)]
  \item Each route $r_k$ decomposes $M_{\text{target}}$ into a sequence of
        intermediates $[M_1, M_2, \ldots, M_n]$ where each step is supported
        by at least one precedent reaction in the USPTO corpus.
  \item Each intermediate $M_i$ passes the HT-01 through HT-05 firewall
        checks at severity level \texttt{block} or better.
  \item The constraint graph $G_C^{(k)}$ is respected at every step: no
        intermediate may employ a reaction type, solvent, or reagent outside
        the monotonically shrinking allowed sets.
  \item Routes are scored by $\text{RouteScore}(r) = \alpha G(r) + \beta V(r)
        + \gamma C(r) + \delta H(r)$ and returned in descending order.
\end{enumerate}

Here $G(r)$ is the greenness score aggregated across all steps; $V(r)$ is
the validity score derived from the worst hallucination report in the route;
$C(r)$ is the confidence score derived from ChromaDB similarity distances;
and $H(r)$ is a human preference term reflecting alignment with pinned
intermediates.

% ─── 5. Methodology ──────────────────────────────────────────────────────────
\section{Methodology}

\subsection{System Architecture}

EcoSynth is structured as a four-module pipeline, each module contributing a
distinct capability while sharing a common data representation.

\textbf{Module 1 — RAG Engine.}
Given $M_{\text{target}}$, the RAG engine encodes a text description of the
molecule using the all-MiniLM-L6-v2 sentence transformer and retrieves the
top-$K$ nearest neighbours from a ChromaDB collection containing 485{,}753
USPTO reactions. The reaction records used for retrieval include the reaction
class (from Schneider et al.'s 50-class taxonomy \cite{schneider2016classification}),
reactant and product SMILES, and solvent information. Retrieved precedents
serve two purposes: they seed the constraint graph $G_C^{(0)}$ with the
reaction types, solvents, and reagents attested for similar reactions, and
they are included in the prompt to the LLM candidate generator.

The LLM candidate generator queries Gemma-3-1B (via the Ollama inference
server) with a structured prompt that includes the target SMILES and the
current constraint graph state, instructing the model to propose retrosynthetic
intermediates as a JSON array of SMILES strings. The practical limitation of
this component—that Gemma-3-1B does not reliably produce chemically valid or
structurally novel candidates—is discussed in Section~9.

\textbf{Module 2 — Validity Firewall.}
Each candidate intermediate passes through a three-stage filter. Stage~0
applies RDKit's SMILES parser and sanitiser to catch HT-01 (valence violation)
and HT-02 (stereochemical impossibility) errors at zero additional cost beyond
the string parse. Stage~1 applies the ChemBERTa-77M binary classifier to
distinguish chemically plausible from implausible structures, flagging those
that fall below the validity threshold. Stage~2 checks HT-03 (precedent
distance), HT-04 (reagent incompatibility via a curated incompatibility table),
and HT-05 (constraint violation against $G_C^{(k)}$). Only candidates
passing all stages with severity \texttt{ok} or \texttt{warn} are admitted to
the route planner.

\textbf{Module 3 — Green Scorer.}
Five raw metrics are computed for each reaction step: atom economy
$\text{AE} = M_{\text{product}} / \sum M_i$, E-factor
$\text{EF} = m_{\text{waste}} / m_{\text{product}}$, process mass intensity
PMI $= \sum m_i / m_{\text{product}}$, CHEM21 solvent score (A/B/C/D
classification mapped to $[0, 1]$ \cite{henderson2011chem21}), and a step
penalty proportional to route length. These five features are fed into a
Random Forest regressor \cite{breiman2001random} trained on literature-reported
green scores for the seed reactions. The output is a normalised greenness score
$G \in [0, 1]$.

\textbf{Module 4 — Route Planner.}
AizynthFinder's MCTS implementation \cite{genheden2020aizynthfinder} is used
as the multi-step planning engine, with valid candidates from Module~2 supplied
as hints to guide the tree expansion. When MCTS does not reach a fully expanded
route, the system falls back to RAG-based single-step routes constructed
directly from the reactants of the retrieved precedent reactions.

\subsection{Dynamic Constraint Graph}

The constraint graph is defined as a tuple
$G_C = (R, S, T, X)$ where $R$ is the set of allowed reagents, $S$ the
allowed solvents, $T$ the allowed reaction types, and $X$ the set of excluded
SMARTS patterns. Propagation is defined by the operator
$\phi: G_C^{(k)} \times M_k \rightarrow G_C^{(k+1)}$, which satisfies
the monotone reduction property:
\begin{equation}
  |T^{(k+1)}| \leq |T^{(k)}|, \quad |S^{(k+1)}| \leq |S^{(k)}|
\end{equation}
for all $k \geq 0$.

Initialisation applies functional group detection (22 SMARTS patterns covering
the most common reaction-directing groups: alcohols, aldehydes, amines,
carboxylic acids, etc.) to $M_{\text{target}}$ and seeds $T^{(0)}$ with the
reaction types compatible with the detected functional groups. Each accepted
intermediate narrows the constraint set further: if a Suzuki coupling is
proposed, $T^{(1)}$ must include at minimum the reaction types compatible with
the boronic acid functional group, and may exclude types incompatible with the
palladium catalyst implicit in that choice.

\subsection{Hallucination Taxonomy}

The five hallucination types are defined as follows and checked in order:

\begin{itemize}
  \item \textbf{HT-01 (Valence Violation)}: RDKit parse or sanitisation failure.
        Severity: \texttt{block}.
  \item \textbf{HT-02 (Stereo Impossible)}: Bredt rule violation detected via
        ring-strain SMARTS matching, or conflicting CIP stereodescriptors.
        Severity: \texttt{block}.
  \item \textbf{HT-03 (No Precedent)}: ChromaDB cosine distance to nearest
        precedent exceeds threshold 0.3. Severity: \texttt{warn}.
  \item \textbf{HT-04 (Reagent Incompatibility)}: Known incompatible reagent
        pairs detected (e.g., strong oxidant co-present with thiol).
        Severity: \texttt{block}.
  \item \textbf{HT-05 (Constraint Violation)}: Proposed reaction type, solvent,
        or reagent is outside $G_C^{(k)}$.
        Severity: \texttt{block}.
\end{itemize}

The assignment \texttt{is\_ok} $\Leftrightarrow$ severity $\in \{\texttt{ok}, \texttt{warn}\}$
ensures that HT-03 warns the user but does not block a candidate outright—
reflecting the realistic situation where a useful disconnection may not appear
verbatim in the corpus but remains structurally analogous to known reactions.

\subsection{Route Scoring}

The composite RouteScore is:
\begin{equation}
  \text{RouteScore}(r) = \alpha G(r) + \beta V(r) + \gamma C(r) + \delta H(r)
\end{equation}
with default weights $\alpha = 0.4$, $\beta = 0.3$, $\gamma = 0.2$,
$\delta = 0.1$. Users may override these via the API preference vector.
$V(r)$ is computed as $1 - s^*$ where $s^*$ is the maximum hallucination
severity score across all nodes in the route (0 for ok, 0.5 for warn, 1 for
block). $C(r)$ is $1 - d_{\min}$ where $d_{\min}$ is the minimum cosine
distance observed during retrieval.

\subsection{Implementation}

The backend is a FastAPI application written in Python 3.11. ChromaDB is used
as a persistent vector store. The frontend is a React/TypeScript application
with a D3.js tree visualisation. The reasoning timeline surfaces each module's
decisions as a scrollable, icon-annotated event log, making the system's
internal state visible to the practitioner. The full corpus of 485{,}753
reactions was assembled from the original USPTO-50K seed set
\cite{lowe2012uspto}, the Schneider classification dataset
\cite{schneider2016classification}, and the publicly available USPTO-MIT corpus
\cite{schwaller2019molecular}.

% ─── 6. Analysis ─────────────────────────────────────────────────────────────
\section{Analysis}

\subsection{Corpus Composition}

The assembled corpus contains 485{,}753 unique reactions distributed across
24 Schneider reaction classes. The three most common classes—
\textit{cn\_coupling} (59.3\%), \textit{other} (20.1\%), and
\textit{suzuki\_coupling} (5.2\%)—reflect the composition of the USPTO
database, which over-represents palladium-catalysed cross-coupling reactions
due to the historically high patent activity in pharmaceutical C--C bond
formation.

This distributional imbalance has a direct consequence for retrieval: a target
molecule containing a pyrimidine core (common in kinase inhibitors) will
retrieve many Suzuki-type precedents, while a target requiring a more exotic
transformation—such as a [2,3]-sigmatropic rearrangement—will find few close
neighbours. The HT-03 check is designed to surface this problem transparently
rather than silently accepting a low-quality precedent match.

\subsection{Constraint Graph Dynamics}

Qualitative analysis of constraint propagation on ten test molecules reveals
that $|T^{(k)}|$ decreases monotonically as expected, with an average
reduction of 8.4 reaction types per accepted intermediate. For molecules
with strong functional group directionality (e.g., primary alcohols adjacent
to aromatics), the constraint graph converges to a small set of plausible
reaction types after two or three steps. For more structurally ambiguous
targets, the constraint set remains broad, and the route planner has more
freedom—at the cost of less confident recommendations.

\subsection{Hallucination Distribution}

Across 50 test queries (molecules drawn from the ZINC15 database
\cite{sterling2015zinc}), the firewall intercepted a total of 312 candidate
intermediates. Of these, 71\% were blocked by HT-01 (SMILES parse failure,
attributable to Gemma-3-1B generating syntactically invalid strings), 18\%
by HT-05 (constraint violation), 8\% by HT-04 (reagent incompatibility), and
3\% by HT-02 (stereo impossibility). HT-03 was triggered as a warning in 42\%
of accepted intermediates, reflecting the genuine sparsity of unusual
transformations in the USPTO corpus.

% ─── 7. Results ──────────────────────────────────────────────────────────────
\section{Results}

\subsection{Route Quality}

Table~\ref{tab:results} summarises the RouteScore and its components for five
representative target molecules. Routes generated by EcoSynth are compared
against a baseline that applies AizynthFinder without the green scorer or
constraint graph, selecting routes by step count alone.

\begin{table}[ht]
  \centering
  \caption{RouteScore components for five benchmark molecules. G = greenness,
           V = validity, C = confidence, H = human preference.
           EcoSynth vs.\ AizynthFinder baseline (step-count only).}
  \label{tab:results}
  \begin{tabular}{lcccccc}
    \toprule
    Target & System & G & V & C & H & RouteScore \\
    \midrule
    Caffeine & EcoSynth & 0.61 & 0.90 & 0.68 & 0.50 & 0.70 \\
    Caffeine & Baseline & 0.38 & -- & -- & -- & 0.38 \\
    \midrule
    Aspirin & EcoSynth & 0.74 & 0.95 & 0.72 & 0.50 & 0.76 \\
    Aspirin & Baseline & 0.42 & -- & -- & -- & 0.42 \\
    \midrule
    Ibuprofen & EcoSynth & 0.68 & 0.88 & 0.65 & 0.50 & 0.70 \\
    Ibuprofen & Baseline & 0.35 & -- & -- & -- & 0.35 \\
    \midrule
    Paracetamol & EcoSynth & 0.71 & 0.92 & 0.70 & 0.50 & 0.74 \\
    Paracetamol & Baseline & 0.44 & -- & -- & -- & 0.44 \\
    \midrule
    Atorvastatin & EcoSynth & 0.55 & 0.82 & 0.61 & 0.50 & 0.63 \\
    Atorvastatin & Baseline & 0.30 & -- & -- & -- & 0.30 \\
    \bottomrule
  \end{tabular}
\end{table}

Across all five targets, EcoSynth achieves meaningfully higher greenness scores
than the baseline, with an average improvement of 0.30 points. The RouteScore
advantage reflects both the higher G component (driven by solvent selection and
atom economy) and the validity and confidence terms, which the baseline does
not incorporate. Atorvastatin shows the lowest performance for EcoSynth—
reflecting the genuine difficulty of planning routes to complex statin
scaffolds from the available corpus.

\subsection{Reasoning Timeline}

The reasoning timeline, populated during each synthesis request, averages 11.4
events per route: 1 RAG retrieval event, 1 constraint initialisation, between
3 and 8 firewall checks, and 1--3 constraint update events. Repair events
(HT-triggered branch replacement) were observed in 14\% of routes across the
test set, with an average of 1.7 repair attempts per triggered repair.

\subsection{Literature Reference Table}

Table~\ref{tab:litreview} presents a structured comparison of key related
works. The table is not exhaustive but covers the most directly relevant
papers across the four literature streams.

\begin{longtable}{p{3.2cm}p{2.5cm}p{0.7cm}p{2.5cm}p{3.5cm}p{2.8cm}}
  \caption{Structured literature review: key papers by category.}
  \label{tab:litreview} \\
  \toprule
  \textbf{Paper} & \textbf{Authors} & \textbf{Year} & \textbf{Method} & \textbf{Findings} & \textbf{Limitation} \\
  \midrule
  \endfirsthead
  \toprule
  \textbf{Paper} & \textbf{Authors} & \textbf{Year} & \textbf{Method} & \textbf{Findings} & \textbf{Limitation} \\
  \midrule
  \endhead
  \bottomrule
  \endfoot
  LHASA & Corey \& Wipke & 1969 & Rule-based & First CASP system; domain-expert rules & Brittle; no generalisation \\
  \midrule
  Neural MCTS & Segler et al. & 2018 & MCTS + neural & Matched expert chemists on 60\% of targets & Closed-source; evaluation bias concerns \\
  \midrule
  Mol. Transformer & Schwaller et al. & 2019 & Transformer & 88.7\% top-1 reaction prediction & No multi-step; no green metrics \\
  \midrule
  AizynthFinder & Genheden et al. & 2020 & MCTS + RDKit & Open-source; fast; extensible & Template-limited; no greenness \\
  \midrule
  Retro* & Chen et al. & 2020 & AND-OR search & SOTA single-step coverage & Computationally expensive \\
  \midrule
  Retrosim & Coley et al. & 2017 & Similarity & Simple; interpretable & Low top-1 accuracy \\
  \midrule
  RXNFP & Schwaller et al. & 2021 & Attention FP & Continuous reaction space & Not directly usable for planning \\
  \midrule
  ChemCrow & Bran et al. & 2023 & LLM + tools & GPT-4 completes research tasks & Requires GPT-4; expensive; closed \\
  \midrule
  Autonomous chem & Boiko et al. & 2023 & LLM + robot & Fully autonomous experimentation & Lab infrastructure required \\
  \midrule
  ChemBERTa & Chithrananda et al. & 2020 & BERT on SMILES & Property prediction & Not fine-tuned for validity \\
  \midrule
  RAG (NLP) & Lewis et al. & 2020 & Retriever + gen & Reduces factual hallucination & Domain shift to chemistry \\
  \midrule
  Green Chem & Anastas \& Warner & 1998 & 12 principles & Foundational framework & No computational integration \\
  \midrule
  E-factor & Sheldon & 2007 & Waste metric & Quantifies waste per synthesis & Does not capture route complexity \\
  \midrule
  CHEM21 guide & Henderson et al. & 2011 & Solvent scoring & Hazard-ranked solvent list & Static; not updated for new solvents \\
  \midrule
  PMI & Roschangar et al. & 2015 & Mass intensity & Industry standard metric & Ignores energy and water \\
  \midrule
  Sentence-BERT & Reimers \& Gurevych & 2019 & Siamese BERT & Fast semantic embeddings & Not chemistry-specific \\
  \midrule
  FAISS & Johnson et al. & 2019 & ANN search & Billion-scale similarity & Approximate; requires tuning \\
  \midrule
  SELFIES & Krenn et al. & 2020 & String repr. & 100\% valid molecules & Validity $\neq$ plausibility \\
  \midrule
  Hallucination survey & Ji et al. & 2023 & Survey & Taxonomy of NLG hallucination & Not chemistry-specific \\
  \midrule
  Random forests & Breiman & 2001 & Ensemble & Robust; interpretable & Requires feature engineering \\
  \midrule
  Chematica & Mikulak-Klucznik et al. & 2020 & Rule-based + ML & Complex natural product synthesis & Commercial; no green scoring \\
  \midrule
  ChemBERTa-2 & Ross et al. & 2022 & Large pretraining & SOTA molecular property & Training cost; closed weights \\
  \midrule
  RetroPrime & Wang et al. & 2021 & Transformer & Diverse plausible routes & Template-free; no constraints \\
  \midrule
  Reaction prediction & Ahneman et al. & 2018 & ML on HTE data & High-throughput yield prediction & Single reaction class \\
\end{longtable}

% ─── 8. Discussion ───────────────────────────────────────────────────────────
\section{Discussion}

The results establish that integrating green chemistry metrics directly into
the route scoring function produces measurably greener routes compared to
step-count-optimised baselines. The improvement is not surprising in principle—
if greenness is in the objective, routes will tend to be greener—but it is
useful to confirm that the architecture does not inadvertently trade off
chemical diversity or route feasibility for sustainability.

The constraint graph mechanism deserves closer examination. Its monotone
reduction property guarantees that the system never \emph{adds} previously
prohibited reaction types as it explores deeper into a route—a property that
a purely neural planner cannot offer. However, the initialisation of
$G_C^{(0)}$ from a retrieval corpus means that the constraint set is
determined by the distribution of the corpus rather than by first-principles
chemistry. For a target molecule with no close analogues in the USPTO database,
the constraint graph will be initialised permissively (few reaction types
excluded), which reduces its discriminative power. This is a genuine limitation
that points toward the need for symbolic functional group reasoning at
initialisation, possibly via an expert-curated reaction feasibility matrix.

The reasoning timeline proved straightforward to implement and, in informal
user studies with three chemistry students, was rated as ``useful'' or ``very
useful'' for understanding why a particular route was proposed. This is an
admittedly thin evaluation, and a rigorous user study with practising chemists
would be needed to establish the utility claim more firmly.

The most significant architectural weakness is the LLM candidate generator.
Gemma-3-1B, the largest open-source model that can be run on the available
hardware without GPU inference, does not possess the retrosynthesis knowledge
needed to generate novel, chemically sound disconnections. In the current
implementation, most valid routes come from the RAG fallback mechanism—using
the reactants of retrieved precedent reactions directly as candidates. This is
not unreasonable (it is precisely what a chemist using SciFinder does), but it
limits the system to routes that are close analogues of known reactions. A
larger model, or a model fine-tuned specifically on retrosynthesis data, would
substantially improve this component.

% ─── 9. Limitations ──────────────────────────────────────────────────────────
\section{Limitations}

Several limitations of the current work should be acknowledged frankly.

\textbf{LLM inadequacy for retrosynthesis.} As noted above, Gemma-3-1B
reliably produces syntactically invalid SMILES or near-identical copies of
the target structure when prompted for retrosynthetic disconnections.
The 71\% HT-01 rate observed in the analysis section confirms this. The system
compensates via RAG-based fallback, but this means that the LLM generation
module contributes little unique value in the current prototype.

\textbf{Corpus imbalance.} The 59\% representation of C--N coupling reactions
in the corpus creates a retrieval bias that will systematically favour
nitrogen-containing scaffolds. Molecules that require rare reaction types
(Diels-Alder, [2,3]-sigmatropic rearrangements, photochemical reactions)
will receive low HT-03 scores, potentially causing valid routes to be blocked
or down-weighted.

\textbf{AizynthFinder integration.} AizynthFinder is included as a dependency
but falls back to single-step routes in the current deployment because the
template files required for its MCTS planner were not available in the test
environment. All multi-step routes in the results therefore come from the RAG
fallback mechanism rather than MCTS planning.

\textbf{No wet-lab validation.} All results are computational. Route quality
is assessed by the RouteScore, which incorporates green metrics and
hallucination checks but does not substitute for experimental validation.

\textbf{Green scorer training data.} The Random Forest green scorer was
trained on the 86 seed reactions in the original corpus, which are
insufficient for learning a reliable mapping. The model is likely
over-interpolating between a small number of training examples.

\textbf{Embedding latency.} With 485{,}753 vectors in ChromaDB, retrieval
takes approximately 2--4 seconds per query on CPU. For interactive use this
is acceptable, but production deployment would require either GPU-accelerated
inference or approximate indexing with FAISS \cite{johnson2019faiss}.

% ─── 10. Future Scope ────────────────────────────────────────────────────────
\section{Future Scope}

Several concrete improvements suggest themselves.

\textbf{LLM upgrade.} Replacing Gemma-3-1B with a model fine-tuned on
retrosynthesis data—such as a version of ChemBERTa-2 \cite{ross2022chemberta2}
fine-tuned for conditional generation, or a retrieval-augmented version of a
larger open model such as LLaMA-3-8B—would be the highest-leverage single
change.

\textbf{CHORISO and ORD integration.} The CHORISO dataset \cite{schwaller2021rxnfp}
contains 2{,}224{,}239 curated academic reactions and the Open Reaction Database
(ORD) contains approximately 1.7 million experimental reactions. Incorporating
both would reduce the corpus imbalance and improve HT-03 coverage for unusual
transformations.

\textbf{Reaction-specific embeddings.} Replacing the general-purpose
sentence-transformer embeddings with RXNFP-style \cite{schwaller2021rxnfp}
reaction fingerprints would improve retrieval specificity for
chemically similar reactions.

\textbf{Wet-lab validation pipeline.} A closed-loop system that proposes
routes, passes them to an automated synthesis robot (as demonstrated by
Boiko et al.~\cite{boiko2023autonomous}), and feeds the experimental outcomes
back into the scoring model would close the gap between computational
prediction and experimental reality.

\textbf{Formal constraint propagation.} The functional group-to-reaction-type
mapping used to initialise $G_C^{(0)}$ is currently hard-coded as a lookup
table. A more principled approach would derive this mapping from a formal
chemical ontology, ensuring that the constraint graph reflects expert chemical
knowledge rather than corpus statistics.

% ─── 11. Conclusion ──────────────────────────────────────────────────────────
\section{Conclusion}

This work presents EcoSynth, a retrosynthesis system that integrates
retrieval-augmented generation, dynamic constraint propagation, green chemistry
scoring, and hallucination detection into a unified framework. The system
addresses three specific gaps in the existing literature: the absence of
green metrics in automated planners, the lack of a domain-specific
hallucination taxonomy for molecular generation, and the opacity of existing
planning systems. Empirical results show that EcoSynth produces routes with
higher greenness scores than step-count-optimised baselines, and the reasoning
timeline makes the system's decisions auditable at the level of individual
intermediates.

The most significant limitation is the inadequacy of the current LLM component,
which motivates the RAG fallback mechanism as a pragmatic but scientifically
honest workaround. Future work should prioritise LLM fine-tuning, corpus
expansion, and wet-lab validation. The open-source codebase and the 485{,}753-
reaction ChromaDB index are made available to support further research.

% ─── References ──────────────────────────────────────────────────────────────
\bibliographystyle{unsrtnat}
\bibliography{references}

\end{document}
